\chapter{Experimental Evaluation}
\label{ch:evaluation}

This chapter evaluates the performance of the proposed document intelligence pipeline across three dimensions: layout detection accuracy, extraction fidelity (specifically regarding complex grids and schematics), and inference latency. The evaluation aims to demonstrate that a specialized ensemble of local models can match or exceed the performance of general-purpose visual APIs when dealing with industrial technical documentation.

\section{Experimental Setup}
\label{sec:experimental_setup}

\subsection{Dataset Composition}
The evaluation utilizes a curated dataset of 250 industrial engineering datasheets and operating manuals from diverse manufacturers (e.g., Siemens, ABB, Schneider Electric). These documents are characterized by high-density tabular data, multi-column reading orders, and embedded CAD-style diagrams. 

Ground truth for extraction was established through manual annotation and subsequent verification by a senior domain expert. For the Knowledge Distillation phase, a separate training set of 1,000 document regions was generated through GPT-4o teacher supervision as described in Chapter \ref{ch:distillation}.

\subsection{Hardware Configuration}
All local inference tasks, including DocLayout-YOLO \cite{doclayout_yolo} detection, Qwen2-VL \cite{qwen25} foundational extraction, and the Two-Pass Markdown Harvest pipeline, were executed on a single workstation equipped with:
\begin{itemize}
    \item \textbf{GPU:} NVIDIA GeForce RTX 3090 (24 GB VRAM)
    \item \textbf{CPU:} AMD Ryzen 9 5950X
    \item \textbf{RAM:} 64 GB DDR4
\end{itemize}
Comparative API baselines were executed using standard cloud-based unstructured document parsing solutions and general-purpose LLM endpoints.

\section{Layout Detection Performance}
\label{sec:layout_results}

The DocLayout-YOLO backbone was evaluated against the general-purpose LayoutLMv3 \cite{layoutlmv3} baseline for region classification and bounding box precision. 

\begin{table}[htbp]
\centering
\begin{tabular}{lccc}
\hline
\textbf{Model} & \textbf{mAP@.5} & \textbf{mAP@.5:.95} & \textbf{Inference (ms)} \\
\hline
LayoutLMv3 (Base) & 0.884 & 0.651 & 420 \\
DocLayout-YOLO (Ours) & \textbf{0.942} & \textbf{0.785} & \textbf{85} \\
\hline
\end{tabular}
\caption{Comparison of layout detection performance. DocLayout-YOLO achieves superior spatial precision and nearly 5x lower latency due to its optimized transformer-detector architecture.}
\label{tab:layout_comparison}
\end{table}

As shown in Table \ref{tab:layout_comparison}, DocLayout-YOLO demonstrates a significant improvement in Mean Average Precision (mAP), particularly in the stricter .5:.95 IoU range. This is primarily attributed to the model's ability to precisely bound overlapping regions such as table captions and figure labels, which are frequently merged by more localized detectors.

\section{Extraction Fidelity and Token Limitation Resolution}
\label{sec:extraction_results}

The primary contribution of this research is the resolution of structural hallucinations and generation token bottlenecks found in ultra-dense parameters tables. We evaluated the extraction accuracy of three configurations on generating structured JSON data:
\begin{enumerate}
    \item \textbf{Commercial Parsing Engine (Docling)}: State-of-the-art visual document processing pipeline mapping directly to structured formats.
    \item \textbf{Single-Pass VLM (GPT-4o)}: Direct extraction via API asking the model to ingest the page and immediately format it to Pydantic JSON schemas.
    \item \textbf{Two-Pass Markdown Harvest (Proposed Pipeline)}: DocLayout-YOLO region masking followed by a pass converting visual data into raw, high-density Markdown via Qwen2-VL, and finally utilizing a structural refiner to format Pydantic schemas.
\end{enumerate}

\begin{table}[htbp]
\centering
\begin{tabular}{lccc}
\hline
\textbf{Metric} & \textbf{Docling Parse} & \textbf{Single-Pass VLM} & \textbf{Two-Pass Harvest (Ours)} \\
\hline
Parameter Recall Rate & 34.5\% & 12.1\% & \textbf{98.5\%} \\
Hallucination Rate & N/A & 42.4\% & \textbf{1.1\%} \\
JSON Validation Pass Rate & 100\% & 8.0\% & \textbf{99.5\%} \\
Average Tokens per Item & N/A & $\sim$20 tokens & \textbf{$\sim$6 tokens} \\
\hline
\end{tabular}
\caption{Extraction fidelity across 500 complex industrial table regions. "Parameter Recall Rate" measures the successful capture of highly specific parameters in deep catalogs without triggering truncation or omissions. "JSON Validation Pass Rate" notes if the file output compiled into valid schemas without syntax errors caused by token crashes.}
\label{tab:extraction_fidelity}
\end{table}

The results in Table \ref{tab:extraction_fidelity} highlight the substantial impact of the Two-Pass architecture. When tasked with extracting 60+ parameters from a single page, Single-Pass VLM extraction fails catastrophically due to maximum context output limits. It truncates JSON outputs yielding an 8.0\% pass rate, or hallucinates schema keys to compensate. By converting dense grids into raw Markdown first, the Proposed Pipeline achieves triple the token efficiency ($\sim$6 tokens per parameter versus 20). This allowed the local ensemble (Qwen2-VL) on consumer hardware to comfortably map all parameters without encountering memory bottlenecks, retaining a 98.5\% recall rate. This demonstrates that specialized pipelining is more effective for domain-specific VRDs than brute-force general-purpose reasoning.

\section{Knowledge Graph and Retrieval Quality}
\label{sec:hkg_results}

To evaluate the impact of the \textit{Hierarchical Knowledge Graph} on downstream RAG applications \cite{rag_lewis}, we measured retrieval precision for "context-dependent" queries (e.g., "What is the supply voltage for the sensor defined in Section 4.2?").

\begin{figure}[htbp]
\centering
\begin{tikzpicture}[
    box/.style={rectangle, draw, thick, align=center, minimum width=3cm, minimum height=0.8cm},
    arrow/.style={->, thick, >=stealth}
]

\node[box, fill=red!10] (flat) {Flat Markdown Chunking\\(Baseline)};
\node[box, right=2cm of flat, fill=green!10] (hkg) {HKG-Based Chunking\\(Ours)};

\node[below=0.5cm of flat] {58\% Retrieval Accuracy};
\node[below=0.5cm of hkg] {94\% Retrieval Accuracy};

\draw[arrow] (flat) -- node[above] {+Breadcrumbs} (hkg);

\end{tikzpicture}
\caption{Impact of the Hierarchical Knowledge Graph on retrieval accuracy. The inclusion of hierarchical breadcrumbs and spatial caption links eliminates ambiguous semantic mapping during vector search.}
\label{fig:retrieval_impact}
\end{figure}

The baseline "Flat Markdown" chunking frequently failed to associate parameters with their correct section headers once the document was segmented. By contrast, the breadcrumb systems in the HKG ensure that the retrieval chunk carries its full lineage, allowing the downstream model to answer complex multi-page queries with 94\% accuracy.

\section{Conclusion on Evaluation}
The empirical results confirm that the proposed pipeline successfully resolves the fundamental structural challenges of industrial VRDs. By combining high-precision layout detection with localized, distilled specialists and graph-native context preservation, the system provides a robust pathway for automating document intelligence on enterprise-scale data.
